\documentclass[aspectratio=169,10pt]{beamer}
\makeatletter\def\input@path{{../../estilo/}{../estilo/}{estilo/}}\makeatother
\usepackage{devmobile}
\modulo[Flutter]{06}{Dart Assíncrono}{Event loop, Future e async/await, tratamento de erros, Stream, StreamController, FutureBuilder/StreamBuilder e Isolates}

\begin{document}
\capa

\begin{frame}{Objetivos do módulo}
  \begin{conceito}[Ao final deste módulo você será capaz de]
    \begin{itemize}
      \item Explicar o \emph{event loop} do Dart e por que a interface não pode ser bloqueada.
      \item Escrever código assíncrono com \cod{Future}, \cod{async} e \cod{await}, tratando erros e tempos limite.
      \item Executar tarefas em paralelo com \cod{Future.wait}.
      \item Produzir e consumir fluxos de dados com \cod{Stream}, \cod{async*} e \cod{StreamController}.
      \item Mostrar dados assíncronos na tela com \cod{FutureBuilder} e \cod{StreamBuilder}.
      \item Levar cálculos pesados para outro \emph{isolate}.
    \end{itemize}
  \end{conceito}
\end{frame}

\sumario

% =============================================================================
%  Fundamentos compartilhados --- Módulo 06: concorrência e assincronia
%  Incluído pelas trilhas Kotlin e Flutter com % =============================================================================
%  Fundamentos compartilhados --- Módulo 06: concorrência e assincronia
%  Incluído pelas trilhas Kotlin e Flutter com % =============================================================================
%  Fundamentos compartilhados --- Módulo 06: concorrência e assincronia
%  Incluído pelas trilhas Kotlin e Flutter com \input{../fundamentos/f06-concorrencia}
% =============================================================================
\section{Fundamentos: concorrência e assincronia}

\begin{frame}{Programas, processos e \emph{threads}}
  \begin{columns}[c]
    \begin{column}{0.47\textwidth}
      \centering
      \begin{tikzpicture}[every node/.style={font=\tiny}]
        \draw[azulpetroleo,line width=0.8pt,rounded corners=4pt,fill=cinzafundo] (0,0) rectangle (5.6,3.2);
        \node[rotulo,anchor=north west] at (0.1,3.15) {\textbf{processo do app}};
        \node[caixaambar,minimum width=5.1cm,minimum height=5mm,font=\tiny] at (2.8,2.45) {memória compartilhada (\emph{heap}): objetos do app};
        \foreach \x/\n in {0.95/{principal (UI)},2.8/{rede},4.65/{cálculo}}{
          \node[caixaescura,minimum width=1.6cm,minimum height=11mm,font=\tiny,align=center] at (\x,1.15) {\emph{thread}\\ \n\\[1pt] \textcolor{ambarclaro}{pilha própria}};
        }
        \node[caixa,minimum width=5.6cm,minimum height=5mm,font=\tiny] (so) at (2.8,-0.55) {sistema operacional: escalonador};
        \foreach \x in {0.7,2.1,3.5,4.9}{\node[caixa,minimum width=1.1cm,minimum height=4mm,font=\tiny] at (\x,-1.3) {núcleo};}
      \end{tikzpicture}
    \end{column}
    \begin{column}{0.51\textwidth}
      \begin{definicao}[Processo e \emph{thread}]
        \textbf{Processo}: programa em execução, com memória própria e isolada dos outros apps. \textbf{\emph{Thread}}: linha de execução dentro do processo; as \emph{threads} dividem a memória do processo, mas cada uma tem sua pilha.
      \end{definicao}
      \begin{pordentro}[O escalonador]
        O sistema operacional distribui as \emph{threads} pelos núcleos e alterna entre elas milhares de vezes por segundo (troca de contexto). Cada \emph{thread} custa memória e tempo de troca: não se cria uma por tarefa.
      \end{pordentro}
      \begin{analogia}
        O processo é a cozinha; as \emph{threads} são os cozinheiros. Eles dividem a despensa, e é aí que começam as brigas pelo mesmo ingrediente.
      \end{analogia}
    \end{column}
  \end{columns}
\end{frame}

\begin{frame}{Concorrência e paralelismo}
  \begin{columns}[c]
    \begin{column}{0.47\textwidth}
      \centering
      \begin{tikzpicture}[x=0.55cm,y=0.5cm,every node/.style={font=\tiny}]
        \node[rotulo,anchor=east] at (-0.2,2) {1 núcleo};
        \foreach \x/\cor in {0/petroleo,1.5/ambar,3/petroleo,4.5/ambar,6/petroleo,7.5/ambar}{\fill[\cor!80] (\x,1.6) rectangle ++(1.4,0.8);}
        \node[rotulo] at (4.5,3.1) {\textbf{concorrência}: as tarefas se alternam};
        \node[rotulo,anchor=east] at (-0.2,0) {núcleo 1};
        \fill[petroleo!80] (0,-0.4) rectangle ++(8.9,0.8);
        \node[rotulo,anchor=east] at (-0.2,-1.1) {núcleo 2};
        \fill[ambar!80] (0,-1.5) rectangle ++(8.9,0.8);
        \node[rotulo] at (4.5,-2.3) {\textbf{paralelismo}: as tarefas rodam juntas};
        \draw[->,cinzamedio] (0,-2.9) -- (9.3,-2.9) node[rotulo,right]{tempo};
        \fill[petroleo!80] (0.4,4) rectangle ++(0.5,0.4); \node[rotulo,anchor=west] at (1,4.2) {tarefa A};
        \fill[ambar!80] (3.4,4) rectangle ++(0.5,0.4); \node[rotulo,anchor=west] at (4,4.2) {tarefa B};
      \end{tikzpicture}
      \raggedright
      \begin{pordentro}[Chips de celular]
        Combinam núcleos de alto desempenho e núcleos econômicos (big.LITTLE): o sistema escolhe onde cada \emph{thread} roda para poupar bateria.
      \end{pordentro}
    \end{column}
    \begin{column}{0.51\textwidth}
      \begin{definicao}[Concorrência]
        Várias tarefas em andamento no mesmo período, que podem se alternar. É um jeito de \textbf{estruturar} o programa.
      \end{definicao}
      \begin{definicao}[Paralelismo]
        Várias tarefas executando \textbf{no mesmo instante}, em núcleos diferentes. É um jeito de \textbf{executar}.
      \end{definicao}
      \begin{historia}[Na frase de Rob Pike (2012)]
        Concorrência é lidar com muitas coisas ao mesmo tempo; paralelismo é fazer muitas coisas ao mesmo tempo.
      \end{historia}
    \end{column}
  \end{columns}
\end{frame}

\begin{frame}{Tarefas de I/O e tarefas de CPU}
  \scriptsize
  \begin{tabularx}{\textwidth}{@{}l L L l l@{}}
    \toprule
    \cabecalho{Tipo} & \cabecalho{O que limita} & \cabecalho{Exemplos} & \cabecalho{Kotlin} & \cabecalho{Dart} \\
    \midrule
    I/O & espera por rede, disco ou banco & requisição HTTP, consulta SQL, ler arquivo & \cod{suspend} + \cod{Dispatchers.IO} & \cod{Future} + \cod{await} \\
    CPU & cálculo no processador & decodificar imagem, JSON gigante, criptografia & \cod{Dispatchers.Default} & \cod{Isolate.run} \\
    Tempo & o relógio & atraso, repetição, tempo limite & \cod{delay}, \cod{withTimeout} & \cod{Future.delayed}, \cod{Timer} \\
    \bottomrule
  \end{tabularx}
  \normalsize
  \vspace{4pt}
  \begin{columns}[T]
    \begin{column}{0.49\textwidth}
      \begin{definicao}[Bloqueante \emph{versus} não bloqueante]
        Uma chamada \textbf{bloqueante} prende a \emph{thread} até terminar. Uma \textbf{não bloqueante} registra o pedido e libera a \emph{thread}; o resultado chega depois, por notificação do sistema.
      \end{definicao}
    \end{column}
    \begin{column}{0.49\textwidth}
      \begin{conceito}[A ideia-chave do módulo]
        \textbf{Esperar não precisa de \emph{thread}}: enquanto a resposta não chega, a mesma \emph{thread} atende outros eventos. Já o \textbf{cálculo pesado precisa de um núcleo} de verdade, fora da \emph{thread} da interface.
      \end{conceito}
    \end{column}
  \end{columns}
\end{frame}

\begin{frame}{O laço de eventos em cada plataforma}
  \begin{columns}[T]
    \begin{column}{0.55\textwidth}
      \begin{definicao}[Laço de eventos (\emph{event loop})]
        Uma \emph{thread} que repete sem parar: pega o próximo evento da fila, executa o código associado até o fim e volta à fila. Toda interface gráfica funciona assim.
      \end{definicao}
      \scriptsize
      \begin{tabularx}{\linewidth}{@{}l L@{}}
        \toprule
        \cabecalho{Plataforma} & \cabecalho{Mecanismo na \emph{thread} principal} \\
        \midrule
        Android & \cod{Looper} + \cod{MessageQueue} + \cod{Handler} \\
        iOS & \cod{RunLoop} e fila principal do GCD \\
        Dart & fila de microtarefas e fila de eventos \\
        Web & laço de eventos do navegador (tarefas e microtarefas) \\
        \bottomrule
      \end{tabularx}
    \end{column}
    \begin{column}{0.43\textwidth}
      \begin{pordentro}[Dart: duas filas]
        Antes de cada evento, o Dart esvazia a fila de \textbf{microtarefas} (continuações curtas de \cod{Future}s). Desde o Flutter 3.29, no Android e no iOS, o código Dart roda na própria \emph{thread} principal do app.
      \end{pordentro}
      \begin{atencao}[Um evento lento atrasa todos]
        Quadros se perdem (\emph{jank}) e, no Android, sem resposta a um toque por cerca de 5\,s aparece o diálogo ANR.
      \end{atencao}
    \end{column}
  \end{columns}
\end{frame}

\begin{frame}{Da pirâmide de \emph{callbacks} ao \cod{async}/\cod{await}}
  \centering
  \begin{tikzpicture}[every node/.style={font=\tiny}]
    \draw[seta] (0,0) -- (13.4,0);
    \foreach \x/\t/\d in {1.2/\emph{callbacks}/{funções chamadas\\ no fim da tarefa},4.4/\emph{Futures} e \emph{Promises}/{objeto que representa\\ o resultado futuro},7.8/{\cod{async}/\cod{await}}/{C\# (2012), Dart (2015),\\ JavaScript (2017)},11.4/concorrência estruturada/{Kotlin (2018):\\ escopos e cancelamento}}{
      \fill[ambar] (\x,0) circle (2pt);
      \node[font=\scriptsize\bfseries,text=azulpetroleo,above=1mm] at (\x,0) {\t};
      \node[rotulo,below=1mm] at (\x,0) {\d};
    }
  \end{tikzpicture}
  \raggedright
  \begin{columns}[T]
    \begin{column}{0.49\textwidth}
      \begin{definicao}[\emph{Future} (ou \emph{Promise})]
        Objeto que representa um resultado que ainda não existe. Começa \textbf{pendente} e termina \textbf{com valor} ou \textbf{com erro}. Quem o recebe pode esperar sem bloquear a \emph{thread}.
      \end{definicao}
      \centering
      \begin{tikzpicture}[x=0.5cm,y=0.45cm,every node/.style={font=\tiny}]
        \node[rotulo,anchor=east] at (-0.2,0.4) {\emph{thread}};
        \fill[petroleo!80] (0,0) rectangle ++(2,0.8); \node[white] at (1,0.4) {trecho 1};
        \fill[cinzaclaro] (2,0) rectangle ++(3,0.8); \node at (3.5,0.4) {outros eventos};
        \fill[petroleo!80] (5,0) rectangle ++(2,0.8); \node[white] at (6,0.4) {trecho 2};
        \fill[cinzaclaro] (7,0) rectangle ++(2.4,0.8); \node at (8.2,0.4) {outros};
        \fill[petroleo!80] (9.4,0) rectangle ++(1.6,0.8); \node[white] at (10.2,0.4) {fim};
        \foreach \x in {2,7}{\draw[ambar,line width=1pt] (\x,-0.2) -- (\x,1.1); \node[rotulo,above] at (\x,1.1) {\cod{await}};}
      \end{tikzpicture}
    \end{column}
    \begin{column}{0.49\textwidth}
      \begin{pordentro}[O que acontece em um \cod{await}]
        A função pausa em um \textbf{ponto de suspensão} e devolve a \emph{thread}; quando o resultado chega, continua de onde parou, com as variáveis locais preservadas. O compilador do Kotlin transforma cada função \cod{suspend} em uma máquina de estados que recebe uma \cod{Continuation}.
      \end{pordentro}
      \begin{dica}
        O ganho é escrever código assíncrono com cara de sequencial: \cod{try}/\cod{catch} e laços funcionam como sempre.
      \end{dica}
    \end{column}
  \end{columns}
\end{frame}

\begin{frame}{Concorrência estruturada}
  \begin{columns}[c]
    \begin{column}{0.45\textwidth}
      \centering
      \begin{tikzpicture}[every node/.style={font=\tiny}]
        \node[caixaescura,minimum width=3.4cm,font=\tiny] (e) at (0,0) {escopo da tela};
        \node[caixa,minimum width=1.6cm,font=\tiny] (a) at (-0.95,-1) {carregar lista};
        \node[caixa,minimum width=1.6cm,font=\tiny] (b) at (0.95,-1) {ouvir sensor};
        \node[caixaambar,minimum width=1.6cm,font=\tiny] (c) at (0.95,-2) {ler histórico};
        \draw[seta] (e) -- (a); \draw[seta] (e) -- (b); \draw[seta] (b) -- (c);
        \draw[pontilhado] (-2,0.5) rectangle (2,-2.5);
        \node[rotulo,text width=4.2cm] at (0,-3) {a tela fecha $\Rightarrow$ o escopo termina $\Rightarrow$ as tarefas filhas são canceladas};
      \end{tikzpicture}
      \begin{historia}
        Assim como a programação estruturada aposentou o \cod{goto}, este modelo aposenta tarefas soltas. A ideia ganhou força em 2018, com um ensaio de Nathaniel J. Smith, e o Kotlin a adotou no mesmo ano.
      \end{historia}
    \end{column}
    \begin{column}{0.53\textwidth}
      \begin{definicao}[Concorrência estruturada]
        Toda tarefa concorrente nasce dentro de um \textbf{escopo} e não vive mais que ele: quando o escopo termina, as tarefas filhas são canceladas.
      \end{definicao}
      \begin{conceito}[Garantias]
        Nada fica rodando depois que a tela fecha; o cancelamento chega aos filhos; erros dos filhos chegam ao pai.
      \end{conceito}
      \begin{atencao}[E no Dart?]
        \cod{Future} e \cod{Stream} não seguem esse modelo: cancele inscrições e temporizadores no \cod{dispose} e confira \cod{mounted} após um \cod{await}.
      \end{atencao}
    \end{column}
  \end{columns}
\end{frame}

\begin{frame}{Condições de corrida e impasses}
  \begin{columns}[T]
    \begin{column}{0.47\textwidth}
      \centering
      \begin{tikzpicture}[x=0.7cm,y=0.6cm,every node/.style={font=\tiny}]
        \tikzset{passo/.style={caixa,minimum height=4.5mm,minimum width=1.2cm,inner sep=2pt,font=\tiny}}
        \node[rotulo,anchor=east] at (-0.1,1) {\emph{thread} A};
        \node[rotulo,anchor=east] at (-0.1,0) {\emph{thread} B};
        \node[passo] at (0.9,1) {lê 5};
        \node[passo] at (2.6,0) {lê 5};
        \node[passo,caixaambar,font=\tiny] at (4.3,1) {grava 6};
        \node[passo,caixaambar,font=\tiny] at (6,0) {grava 6};
        \draw[->,cinzamedio] (0,-0.7) -- (7,-0.7) node[rotulo,right]{t};
        \node[rotulo,text=vermelhoerro,text width=4.8cm] at (3.4,-1.3) {duas somas de 1 a partir de 5 deveriam dar 7: uma atualização se perdeu};
      \end{tikzpicture}
      \begin{definicao}[Condição de corrida]
        Erro em que o resultado depende da ordem imprevisível em que tarefas concorrentes acessam um dado compartilhado.
      \end{definicao}
      \begin{atencao}[Impasse (\emph{deadlock})]
        Duas tarefas, cada uma segurando um recurso e esperando o da outra: ninguém avança. Adquira travas sempre na mesma ordem.
      \end{atencao}
    \end{column}
    \begin{column}{0.51\textwidth}
      \scriptsize
      \begin{tabularx}{\linewidth}{@{}l L@{}}
        \toprule
        \cabecalho{Proteção} & \cabecalho{Como funciona} \\
        \midrule
        Exclusão mútua & uma tarefa por vez na região crítica (\cod{Mutex}) \\
        Operação atômica & ler e gravar em um passo indivisível (\cod{AtomicInteger}, \cod{StateFlow.update}) \\
        Confinamento & só uma \emph{thread} toca o dado (a principal, por exemplo) \\
        Imutabilidade & ninguém altera: cria-se uma cópia nova \\
        Mensagens & cada tarefa tem seus dados e troca cópias \\
        \bottomrule
      \end{tabularx}
      \normalsize
      \begin{pordentro}[Corrida lógica com uma \emph{thread} só]
        Entre dois \cod{await}, outro evento pode mudar o estado. Duas buscas disparadas em sequência podem responder fora de ordem, e a mais antiga sobrescrever a mais nova.
      \end{pordentro}
    \end{column}
  \end{columns}
\end{frame}

\begin{frame}{Memória compartilhada ou troca de mensagens}
  \begin{columns}[c]
    \begin{column}{0.43\textwidth}
      \centering
      \begin{tikzpicture}[every node/.style={font=\tiny}]
        \node[caixa,minimum width=1.3cm,font=\tiny] (a) at (0,0) {\emph{thread} A};
        \node[caixa,minimum width=1.3cm,font=\tiny] (b) at (2.6,0) {\emph{thread} B};
        \node[caixaambar,minimum width=1.6cm,font=\tiny] (o) at (1.3,-1.2) {objeto + trava};
        \draw[seta] (a) -- (o); \draw[seta] (b) -- (o);
        \node[rotulo] at (1.3,0.6) {\textbf{memória compartilhada}};
        \node[caixaescura,minimum width=1.3cm,minimum height=9mm,font=\tiny] (c) at (0,-3) {isolate A\\ memória A};
        \node[caixaescura,minimum width=1.3cm,minimum height=9mm,font=\tiny] (d) at (2.6,-3) {isolate B\\ memória B};
        \draw[setaambar] ([yshift=1.5mm]c.east) -- node[rotulo,above]{mensagem} ([yshift=1.5mm]d.west);
        \draw[setaambar] ([yshift=-1.5mm]d.west) -- node[rotulo,below]{resposta} ([yshift=-1.5mm]c.east);
        \node[rotulo] at (1.3,-2.1) {\textbf{troca de mensagens}};
      \end{tikzpicture}
    \end{column}
    \begin{column}{0.55\textwidth}
      \scriptsize
      \begin{tabularx}{\linewidth}{@{}l L L@{}}
        \toprule
         & \cabecalho{Memória compartilhada} & \cabecalho{Troca de mensagens} \\
        \midrule
        Exemplos & \emph{threads} da JVM (Kotlin, Java) & \emph{isolates} do Dart, atores, Web Workers \\
        Comunicação & ler e gravar o mesmo objeto & enviar cópias de dados \\
        Risco & corridas e impasses & custo de copiar dados grandes \\
        Proteção & travas, atômicos, confinamento & o próprio modelo isola \\
        \bottomrule
      \end{tabularx}
      \normalsize
      \begin{analogia}
        Um quadro branco compartilhado exige regras para ninguém apagar o do colega. Colegas que trocam bilhetes nunca brigam pelo quadro, mas gastam tempo escrevendo.
      \end{analogia}
    \end{column}
  \end{columns}
\end{frame}

\begin{frame}{Fluxos de dados: frios, quentes e contrapressão}
  \begin{columns}[T]
    \begin{column}{0.55\textwidth}
      \begin{definicao}[Fluxo (\emph{stream})]
        Sequência de valores que chegam ao longo do tempo: leituras de sensor, mensagens, mudanças no banco. Segue o padrão \textbf{Observador}: quem produz notifica quem se inscreveu.
      \end{definicao}
      \scriptsize
      \begin{tabularx}{\linewidth}{@{}l L l l@{}}
        \toprule
        \cabecalho{Tipo} & \cabecalho{Comportamento} & \cabecalho{Kotlin} & \cabecalho{Dart} \\
        \midrule
        Frio & o produtor começa a cada inscrição & \cod{Flow} & \cod{Stream} (\cod{async*}) \\
        Quente & existe sem observadores; quem chega tarde perde o passado & \cod{SharedFlow} & \cod{Stream} \emph{broadcast} \\
        Estado & quente e guarda o valor atual & \cod{StateFlow} & \cod{ValueNotifier} \\
        \bottomrule
      \end{tabularx}
    \end{column}
    \begin{column}{0.43\textwidth}
      \begin{definicao}[Contrapressão (\emph{backpressure})]
        O que fazer quando o produtor é mais rápido que o consumidor: acumular (\emph{buffer}), descartar, manter só o mais recente (\emph{conflate}) ou amostrar (\emph{debounce}, \emph{sample}).
      \end{definicao}
      \begin{exemplo}[No LabMonitor]
        Um sensor envia 100 leituras por segundo e a tela desenha 60 quadros: basta mostrar sempre a leitura mais recente.
      \end{exemplo}
    \end{column}
  \end{columns}
\end{frame}

% =============================================================================
\section{Fundamentos: concorrência e assincronia}

\begin{frame}{Programas, processos e \emph{threads}}
  \begin{columns}[c]
    \begin{column}{0.47\textwidth}
      \centering
      \begin{tikzpicture}[every node/.style={font=\tiny}]
        \draw[azulpetroleo,line width=0.8pt,rounded corners=4pt,fill=cinzafundo] (0,0) rectangle (5.6,3.2);
        \node[rotulo,anchor=north west] at (0.1,3.15) {\textbf{processo do app}};
        \node[caixaambar,minimum width=5.1cm,minimum height=5mm,font=\tiny] at (2.8,2.45) {memória compartilhada (\emph{heap}): objetos do app};
        \foreach \x/\n in {0.95/{principal (UI)},2.8/{rede},4.65/{cálculo}}{
          \node[caixaescura,minimum width=1.6cm,minimum height=11mm,font=\tiny,align=center] at (\x,1.15) {\emph{thread}\\ \n\\[1pt] \textcolor{ambarclaro}{pilha própria}};
        }
        \node[caixa,minimum width=5.6cm,minimum height=5mm,font=\tiny] (so) at (2.8,-0.55) {sistema operacional: escalonador};
        \foreach \x in {0.7,2.1,3.5,4.9}{\node[caixa,minimum width=1.1cm,minimum height=4mm,font=\tiny] at (\x,-1.3) {núcleo};}
      \end{tikzpicture}
    \end{column}
    \begin{column}{0.51\textwidth}
      \begin{definicao}[Processo e \emph{thread}]
        \textbf{Processo}: programa em execução, com memória própria e isolada dos outros apps. \textbf{\emph{Thread}}: linha de execução dentro do processo; as \emph{threads} dividem a memória do processo, mas cada uma tem sua pilha.
      \end{definicao}
      \begin{pordentro}[O escalonador]
        O sistema operacional distribui as \emph{threads} pelos núcleos e alterna entre elas milhares de vezes por segundo (troca de contexto). Cada \emph{thread} custa memória e tempo de troca: não se cria uma por tarefa.
      \end{pordentro}
      \begin{analogia}
        O processo é a cozinha; as \emph{threads} são os cozinheiros. Eles dividem a despensa, e é aí que começam as brigas pelo mesmo ingrediente.
      \end{analogia}
    \end{column}
  \end{columns}
\end{frame}

\begin{frame}{Concorrência e paralelismo}
  \begin{columns}[c]
    \begin{column}{0.47\textwidth}
      \centering
      \begin{tikzpicture}[x=0.55cm,y=0.5cm,every node/.style={font=\tiny}]
        \node[rotulo,anchor=east] at (-0.2,2) {1 núcleo};
        \foreach \x/\cor in {0/petroleo,1.5/ambar,3/petroleo,4.5/ambar,6/petroleo,7.5/ambar}{\fill[\cor!80] (\x,1.6) rectangle ++(1.4,0.8);}
        \node[rotulo] at (4.5,3.1) {\textbf{concorrência}: as tarefas se alternam};
        \node[rotulo,anchor=east] at (-0.2,0) {núcleo 1};
        \fill[petroleo!80] (0,-0.4) rectangle ++(8.9,0.8);
        \node[rotulo,anchor=east] at (-0.2,-1.1) {núcleo 2};
        \fill[ambar!80] (0,-1.5) rectangle ++(8.9,0.8);
        \node[rotulo] at (4.5,-2.3) {\textbf{paralelismo}: as tarefas rodam juntas};
        \draw[->,cinzamedio] (0,-2.9) -- (9.3,-2.9) node[rotulo,right]{tempo};
        \fill[petroleo!80] (0.4,4) rectangle ++(0.5,0.4); \node[rotulo,anchor=west] at (1,4.2) {tarefa A};
        \fill[ambar!80] (3.4,4) rectangle ++(0.5,0.4); \node[rotulo,anchor=west] at (4,4.2) {tarefa B};
      \end{tikzpicture}
      \raggedright
      \begin{pordentro}[Chips de celular]
        Combinam núcleos de alto desempenho e núcleos econômicos (big.LITTLE): o sistema escolhe onde cada \emph{thread} roda para poupar bateria.
      \end{pordentro}
    \end{column}
    \begin{column}{0.51\textwidth}
      \begin{definicao}[Concorrência]
        Várias tarefas em andamento no mesmo período, que podem se alternar. É um jeito de \textbf{estruturar} o programa.
      \end{definicao}
      \begin{definicao}[Paralelismo]
        Várias tarefas executando \textbf{no mesmo instante}, em núcleos diferentes. É um jeito de \textbf{executar}.
      \end{definicao}
      \begin{historia}[Na frase de Rob Pike (2012)]
        Concorrência é lidar com muitas coisas ao mesmo tempo; paralelismo é fazer muitas coisas ao mesmo tempo.
      \end{historia}
    \end{column}
  \end{columns}
\end{frame}

\begin{frame}{Tarefas de I/O e tarefas de CPU}
  \scriptsize
  \begin{tabularx}{\textwidth}{@{}l L L l l@{}}
    \toprule
    \cabecalho{Tipo} & \cabecalho{O que limita} & \cabecalho{Exemplos} & \cabecalho{Kotlin} & \cabecalho{Dart} \\
    \midrule
    I/O & espera por rede, disco ou banco & requisição HTTP, consulta SQL, ler arquivo & \cod{suspend} + \cod{Dispatchers.IO} & \cod{Future} + \cod{await} \\
    CPU & cálculo no processador & decodificar imagem, JSON gigante, criptografia & \cod{Dispatchers.Default} & \cod{Isolate.run} \\
    Tempo & o relógio & atraso, repetição, tempo limite & \cod{delay}, \cod{withTimeout} & \cod{Future.delayed}, \cod{Timer} \\
    \bottomrule
  \end{tabularx}
  \normalsize
  \vspace{4pt}
  \begin{columns}[T]
    \begin{column}{0.49\textwidth}
      \begin{definicao}[Bloqueante \emph{versus} não bloqueante]
        Uma chamada \textbf{bloqueante} prende a \emph{thread} até terminar. Uma \textbf{não bloqueante} registra o pedido e libera a \emph{thread}; o resultado chega depois, por notificação do sistema.
      \end{definicao}
    \end{column}
    \begin{column}{0.49\textwidth}
      \begin{conceito}[A ideia-chave do módulo]
        \textbf{Esperar não precisa de \emph{thread}}: enquanto a resposta não chega, a mesma \emph{thread} atende outros eventos. Já o \textbf{cálculo pesado precisa de um núcleo} de verdade, fora da \emph{thread} da interface.
      \end{conceito}
    \end{column}
  \end{columns}
\end{frame}

\begin{frame}{O laço de eventos em cada plataforma}
  \begin{columns}[T]
    \begin{column}{0.55\textwidth}
      \begin{definicao}[Laço de eventos (\emph{event loop})]
        Uma \emph{thread} que repete sem parar: pega o próximo evento da fila, executa o código associado até o fim e volta à fila. Toda interface gráfica funciona assim.
      \end{definicao}
      \scriptsize
      \begin{tabularx}{\linewidth}{@{}l L@{}}
        \toprule
        \cabecalho{Plataforma} & \cabecalho{Mecanismo na \emph{thread} principal} \\
        \midrule
        Android & \cod{Looper} + \cod{MessageQueue} + \cod{Handler} \\
        iOS & \cod{RunLoop} e fila principal do GCD \\
        Dart & fila de microtarefas e fila de eventos \\
        Web & laço de eventos do navegador (tarefas e microtarefas) \\
        \bottomrule
      \end{tabularx}
    \end{column}
    \begin{column}{0.43\textwidth}
      \begin{pordentro}[Dart: duas filas]
        Antes de cada evento, o Dart esvazia a fila de \textbf{microtarefas} (continuações curtas de \cod{Future}s). Desde o Flutter 3.29, no Android e no iOS, o código Dart roda na própria \emph{thread} principal do app.
      \end{pordentro}
      \begin{atencao}[Um evento lento atrasa todos]
        Quadros se perdem (\emph{jank}) e, no Android, sem resposta a um toque por cerca de 5\,s aparece o diálogo ANR.
      \end{atencao}
    \end{column}
  \end{columns}
\end{frame}

\begin{frame}{Da pirâmide de \emph{callbacks} ao \cod{async}/\cod{await}}
  \centering
  \begin{tikzpicture}[every node/.style={font=\tiny}]
    \draw[seta] (0,0) -- (13.4,0);
    \foreach \x/\t/\d in {1.2/\emph{callbacks}/{funções chamadas\\ no fim da tarefa},4.4/\emph{Futures} e \emph{Promises}/{objeto que representa\\ o resultado futuro},7.8/{\cod{async}/\cod{await}}/{C\# (2012), Dart (2015),\\ JavaScript (2017)},11.4/concorrência estruturada/{Kotlin (2018):\\ escopos e cancelamento}}{
      \fill[ambar] (\x,0) circle (2pt);
      \node[font=\scriptsize\bfseries,text=azulpetroleo,above=1mm] at (\x,0) {\t};
      \node[rotulo,below=1mm] at (\x,0) {\d};
    }
  \end{tikzpicture}
  \raggedright
  \begin{columns}[T]
    \begin{column}{0.49\textwidth}
      \begin{definicao}[\emph{Future} (ou \emph{Promise})]
        Objeto que representa um resultado que ainda não existe. Começa \textbf{pendente} e termina \textbf{com valor} ou \textbf{com erro}. Quem o recebe pode esperar sem bloquear a \emph{thread}.
      \end{definicao}
      \centering
      \begin{tikzpicture}[x=0.5cm,y=0.45cm,every node/.style={font=\tiny}]
        \node[rotulo,anchor=east] at (-0.2,0.4) {\emph{thread}};
        \fill[petroleo!80] (0,0) rectangle ++(2,0.8); \node[white] at (1,0.4) {trecho 1};
        \fill[cinzaclaro] (2,0) rectangle ++(3,0.8); \node at (3.5,0.4) {outros eventos};
        \fill[petroleo!80] (5,0) rectangle ++(2,0.8); \node[white] at (6,0.4) {trecho 2};
        \fill[cinzaclaro] (7,0) rectangle ++(2.4,0.8); \node at (8.2,0.4) {outros};
        \fill[petroleo!80] (9.4,0) rectangle ++(1.6,0.8); \node[white] at (10.2,0.4) {fim};
        \foreach \x in {2,7}{\draw[ambar,line width=1pt] (\x,-0.2) -- (\x,1.1); \node[rotulo,above] at (\x,1.1) {\cod{await}};}
      \end{tikzpicture}
    \end{column}
    \begin{column}{0.49\textwidth}
      \begin{pordentro}[O que acontece em um \cod{await}]
        A função pausa em um \textbf{ponto de suspensão} e devolve a \emph{thread}; quando o resultado chega, continua de onde parou, com as variáveis locais preservadas. O compilador do Kotlin transforma cada função \cod{suspend} em uma máquina de estados que recebe uma \cod{Continuation}.
      \end{pordentro}
      \begin{dica}
        O ganho é escrever código assíncrono com cara de sequencial: \cod{try}/\cod{catch} e laços funcionam como sempre.
      \end{dica}
    \end{column}
  \end{columns}
\end{frame}

\begin{frame}{Concorrência estruturada}
  \begin{columns}[c]
    \begin{column}{0.45\textwidth}
      \centering
      \begin{tikzpicture}[every node/.style={font=\tiny}]
        \node[caixaescura,minimum width=3.4cm,font=\tiny] (e) at (0,0) {escopo da tela};
        \node[caixa,minimum width=1.6cm,font=\tiny] (a) at (-0.95,-1) {carregar lista};
        \node[caixa,minimum width=1.6cm,font=\tiny] (b) at (0.95,-1) {ouvir sensor};
        \node[caixaambar,minimum width=1.6cm,font=\tiny] (c) at (0.95,-2) {ler histórico};
        \draw[seta] (e) -- (a); \draw[seta] (e) -- (b); \draw[seta] (b) -- (c);
        \draw[pontilhado] (-2,0.5) rectangle (2,-2.5);
        \node[rotulo,text width=4.2cm] at (0,-3) {a tela fecha $\Rightarrow$ o escopo termina $\Rightarrow$ as tarefas filhas são canceladas};
      \end{tikzpicture}
      \begin{historia}
        Assim como a programação estruturada aposentou o \cod{goto}, este modelo aposenta tarefas soltas. A ideia ganhou força em 2018, com um ensaio de Nathaniel J. Smith, e o Kotlin a adotou no mesmo ano.
      \end{historia}
    \end{column}
    \begin{column}{0.53\textwidth}
      \begin{definicao}[Concorrência estruturada]
        Toda tarefa concorrente nasce dentro de um \textbf{escopo} e não vive mais que ele: quando o escopo termina, as tarefas filhas são canceladas.
      \end{definicao}
      \begin{conceito}[Garantias]
        Nada fica rodando depois que a tela fecha; o cancelamento chega aos filhos; erros dos filhos chegam ao pai.
      \end{conceito}
      \begin{atencao}[E no Dart?]
        \cod{Future} e \cod{Stream} não seguem esse modelo: cancele inscrições e temporizadores no \cod{dispose} e confira \cod{mounted} após um \cod{await}.
      \end{atencao}
    \end{column}
  \end{columns}
\end{frame}

\begin{frame}{Condições de corrida e impasses}
  \begin{columns}[T]
    \begin{column}{0.47\textwidth}
      \centering
      \begin{tikzpicture}[x=0.7cm,y=0.6cm,every node/.style={font=\tiny}]
        \tikzset{passo/.style={caixa,minimum height=4.5mm,minimum width=1.2cm,inner sep=2pt,font=\tiny}}
        \node[rotulo,anchor=east] at (-0.1,1) {\emph{thread} A};
        \node[rotulo,anchor=east] at (-0.1,0) {\emph{thread} B};
        \node[passo] at (0.9,1) {lê 5};
        \node[passo] at (2.6,0) {lê 5};
        \node[passo,caixaambar,font=\tiny] at (4.3,1) {grava 6};
        \node[passo,caixaambar,font=\tiny] at (6,0) {grava 6};
        \draw[->,cinzamedio] (0,-0.7) -- (7,-0.7) node[rotulo,right]{t};
        \node[rotulo,text=vermelhoerro,text width=4.8cm] at (3.4,-1.3) {duas somas de 1 a partir de 5 deveriam dar 7: uma atualização se perdeu};
      \end{tikzpicture}
      \begin{definicao}[Condição de corrida]
        Erro em que o resultado depende da ordem imprevisível em que tarefas concorrentes acessam um dado compartilhado.
      \end{definicao}
      \begin{atencao}[Impasse (\emph{deadlock})]
        Duas tarefas, cada uma segurando um recurso e esperando o da outra: ninguém avança. Adquira travas sempre na mesma ordem.
      \end{atencao}
    \end{column}
    \begin{column}{0.51\textwidth}
      \scriptsize
      \begin{tabularx}{\linewidth}{@{}l L@{}}
        \toprule
        \cabecalho{Proteção} & \cabecalho{Como funciona} \\
        \midrule
        Exclusão mútua & uma tarefa por vez na região crítica (\cod{Mutex}) \\
        Operação atômica & ler e gravar em um passo indivisível (\cod{AtomicInteger}, \cod{StateFlow.update}) \\
        Confinamento & só uma \emph{thread} toca o dado (a principal, por exemplo) \\
        Imutabilidade & ninguém altera: cria-se uma cópia nova \\
        Mensagens & cada tarefa tem seus dados e troca cópias \\
        \bottomrule
      \end{tabularx}
      \normalsize
      \begin{pordentro}[Corrida lógica com uma \emph{thread} só]
        Entre dois \cod{await}, outro evento pode mudar o estado. Duas buscas disparadas em sequência podem responder fora de ordem, e a mais antiga sobrescrever a mais nova.
      \end{pordentro}
    \end{column}
  \end{columns}
\end{frame}

\begin{frame}{Memória compartilhada ou troca de mensagens}
  \begin{columns}[c]
    \begin{column}{0.43\textwidth}
      \centering
      \begin{tikzpicture}[every node/.style={font=\tiny}]
        \node[caixa,minimum width=1.3cm,font=\tiny] (a) at (0,0) {\emph{thread} A};
        \node[caixa,minimum width=1.3cm,font=\tiny] (b) at (2.6,0) {\emph{thread} B};
        \node[caixaambar,minimum width=1.6cm,font=\tiny] (o) at (1.3,-1.2) {objeto + trava};
        \draw[seta] (a) -- (o); \draw[seta] (b) -- (o);
        \node[rotulo] at (1.3,0.6) {\textbf{memória compartilhada}};
        \node[caixaescura,minimum width=1.3cm,minimum height=9mm,font=\tiny] (c) at (0,-3) {isolate A\\ memória A};
        \node[caixaescura,minimum width=1.3cm,minimum height=9mm,font=\tiny] (d) at (2.6,-3) {isolate B\\ memória B};
        \draw[setaambar] ([yshift=1.5mm]c.east) -- node[rotulo,above]{mensagem} ([yshift=1.5mm]d.west);
        \draw[setaambar] ([yshift=-1.5mm]d.west) -- node[rotulo,below]{resposta} ([yshift=-1.5mm]c.east);
        \node[rotulo] at (1.3,-2.1) {\textbf{troca de mensagens}};
      \end{tikzpicture}
    \end{column}
    \begin{column}{0.55\textwidth}
      \scriptsize
      \begin{tabularx}{\linewidth}{@{}l L L@{}}
        \toprule
         & \cabecalho{Memória compartilhada} & \cabecalho{Troca de mensagens} \\
        \midrule
        Exemplos & \emph{threads} da JVM (Kotlin, Java) & \emph{isolates} do Dart, atores, Web Workers \\
        Comunicação & ler e gravar o mesmo objeto & enviar cópias de dados \\
        Risco & corridas e impasses & custo de copiar dados grandes \\
        Proteção & travas, atômicos, confinamento & o próprio modelo isola \\
        \bottomrule
      \end{tabularx}
      \normalsize
      \begin{analogia}
        Um quadro branco compartilhado exige regras para ninguém apagar o do colega. Colegas que trocam bilhetes nunca brigam pelo quadro, mas gastam tempo escrevendo.
      \end{analogia}
    \end{column}
  \end{columns}
\end{frame}

\begin{frame}{Fluxos de dados: frios, quentes e contrapressão}
  \begin{columns}[T]
    \begin{column}{0.55\textwidth}
      \begin{definicao}[Fluxo (\emph{stream})]
        Sequência de valores que chegam ao longo do tempo: leituras de sensor, mensagens, mudanças no banco. Segue o padrão \textbf{Observador}: quem produz notifica quem se inscreveu.
      \end{definicao}
      \scriptsize
      \begin{tabularx}{\linewidth}{@{}l L l l@{}}
        \toprule
        \cabecalho{Tipo} & \cabecalho{Comportamento} & \cabecalho{Kotlin} & \cabecalho{Dart} \\
        \midrule
        Frio & o produtor começa a cada inscrição & \cod{Flow} & \cod{Stream} (\cod{async*}) \\
        Quente & existe sem observadores; quem chega tarde perde o passado & \cod{SharedFlow} & \cod{Stream} \emph{broadcast} \\
        Estado & quente e guarda o valor atual & \cod{StateFlow} & \cod{ValueNotifier} \\
        \bottomrule
      \end{tabularx}
    \end{column}
    \begin{column}{0.43\textwidth}
      \begin{definicao}[Contrapressão (\emph{backpressure})]
        O que fazer quando o produtor é mais rápido que o consumidor: acumular (\emph{buffer}), descartar, manter só o mais recente (\emph{conflate}) ou amostrar (\emph{debounce}, \emph{sample}).
      \end{definicao}
      \begin{exemplo}[No LabMonitor]
        Um sensor envia 100 leituras por segundo e a tela desenha 60 quadros: basta mostrar sempre a leitura mais recente.
      \end{exemplo}
    \end{column}
  \end{columns}
\end{frame}

% =============================================================================
\section{Fundamentos: concorrência e assincronia}

\begin{frame}{Programas, processos e \emph{threads}}
  \begin{columns}[c]
    \begin{column}{0.47\textwidth}
      \centering
      \begin{tikzpicture}[every node/.style={font=\tiny}]
        \draw[azulpetroleo,line width=0.8pt,rounded corners=4pt,fill=cinzafundo] (0,0) rectangle (5.6,3.2);
        \node[rotulo,anchor=north west] at (0.1,3.15) {\textbf{processo do app}};
        \node[caixaambar,minimum width=5.1cm,minimum height=5mm,font=\tiny] at (2.8,2.45) {memória compartilhada (\emph{heap}): objetos do app};
        \foreach \x/\n in {0.95/{principal (UI)},2.8/{rede},4.65/{cálculo}}{
          \node[caixaescura,minimum width=1.6cm,minimum height=11mm,font=\tiny,align=center] at (\x,1.15) {\emph{thread}\\ \n\\[1pt] \textcolor{ambarclaro}{pilha própria}};
        }
        \node[caixa,minimum width=5.6cm,minimum height=5mm,font=\tiny] (so) at (2.8,-0.55) {sistema operacional: escalonador};
        \foreach \x in {0.7,2.1,3.5,4.9}{\node[caixa,minimum width=1.1cm,minimum height=4mm,font=\tiny] at (\x,-1.3) {núcleo};}
      \end{tikzpicture}
    \end{column}
    \begin{column}{0.51\textwidth}
      \begin{definicao}[Processo e \emph{thread}]
        \textbf{Processo}: programa em execução, com memória própria e isolada dos outros apps. \textbf{\emph{Thread}}: linha de execução dentro do processo; as \emph{threads} dividem a memória do processo, mas cada uma tem sua pilha.
      \end{definicao}
      \begin{pordentro}[O escalonador]
        O sistema operacional distribui as \emph{threads} pelos núcleos e alterna entre elas milhares de vezes por segundo (troca de contexto). Cada \emph{thread} custa memória e tempo de troca: não se cria uma por tarefa.
      \end{pordentro}
      \begin{analogia}
        O processo é a cozinha; as \emph{threads} são os cozinheiros. Eles dividem a despensa, e é aí que começam as brigas pelo mesmo ingrediente.
      \end{analogia}
    \end{column}
  \end{columns}
\end{frame}

\begin{frame}{Concorrência e paralelismo}
  \begin{columns}[c]
    \begin{column}{0.47\textwidth}
      \centering
      \begin{tikzpicture}[x=0.55cm,y=0.5cm,every node/.style={font=\tiny}]
        \node[rotulo,anchor=east] at (-0.2,2) {1 núcleo};
        \foreach \x/\cor in {0/petroleo,1.5/ambar,3/petroleo,4.5/ambar,6/petroleo,7.5/ambar}{\fill[\cor!80] (\x,1.6) rectangle ++(1.4,0.8);}
        \node[rotulo] at (4.5,3.1) {\textbf{concorrência}: as tarefas se alternam};
        \node[rotulo,anchor=east] at (-0.2,0) {núcleo 1};
        \fill[petroleo!80] (0,-0.4) rectangle ++(8.9,0.8);
        \node[rotulo,anchor=east] at (-0.2,-1.1) {núcleo 2};
        \fill[ambar!80] (0,-1.5) rectangle ++(8.9,0.8);
        \node[rotulo] at (4.5,-2.3) {\textbf{paralelismo}: as tarefas rodam juntas};
        \draw[->,cinzamedio] (0,-2.9) -- (9.3,-2.9) node[rotulo,right]{tempo};
        \fill[petroleo!80] (0.4,4) rectangle ++(0.5,0.4); \node[rotulo,anchor=west] at (1,4.2) {tarefa A};
        \fill[ambar!80] (3.4,4) rectangle ++(0.5,0.4); \node[rotulo,anchor=west] at (4,4.2) {tarefa B};
      \end{tikzpicture}
      \raggedright
      \begin{pordentro}[Chips de celular]
        Combinam núcleos de alto desempenho e núcleos econômicos (big.LITTLE): o sistema escolhe onde cada \emph{thread} roda para poupar bateria.
      \end{pordentro}
    \end{column}
    \begin{column}{0.51\textwidth}
      \begin{definicao}[Concorrência]
        Várias tarefas em andamento no mesmo período, que podem se alternar. É um jeito de \textbf{estruturar} o programa.
      \end{definicao}
      \begin{definicao}[Paralelismo]
        Várias tarefas executando \textbf{no mesmo instante}, em núcleos diferentes. É um jeito de \textbf{executar}.
      \end{definicao}
      \begin{historia}[Na frase de Rob Pike (2012)]
        Concorrência é lidar com muitas coisas ao mesmo tempo; paralelismo é fazer muitas coisas ao mesmo tempo.
      \end{historia}
    \end{column}
  \end{columns}
\end{frame}

\begin{frame}{Tarefas de I/O e tarefas de CPU}
  \scriptsize
  \begin{tabularx}{\textwidth}{@{}l L L l l@{}}
    \toprule
    \cabecalho{Tipo} & \cabecalho{O que limita} & \cabecalho{Exemplos} & \cabecalho{Kotlin} & \cabecalho{Dart} \\
    \midrule
    I/O & espera por rede, disco ou banco & requisição HTTP, consulta SQL, ler arquivo & \cod{suspend} + \cod{Dispatchers.IO} & \cod{Future} + \cod{await} \\
    CPU & cálculo no processador & decodificar imagem, JSON gigante, criptografia & \cod{Dispatchers.Default} & \cod{Isolate.run} \\
    Tempo & o relógio & atraso, repetição, tempo limite & \cod{delay}, \cod{withTimeout} & \cod{Future.delayed}, \cod{Timer} \\
    \bottomrule
  \end{tabularx}
  \normalsize
  \vspace{4pt}
  \begin{columns}[T]
    \begin{column}{0.49\textwidth}
      \begin{definicao}[Bloqueante \emph{versus} não bloqueante]
        Uma chamada \textbf{bloqueante} prende a \emph{thread} até terminar. Uma \textbf{não bloqueante} registra o pedido e libera a \emph{thread}; o resultado chega depois, por notificação do sistema.
      \end{definicao}
    \end{column}
    \begin{column}{0.49\textwidth}
      \begin{conceito}[A ideia-chave do módulo]
        \textbf{Esperar não precisa de \emph{thread}}: enquanto a resposta não chega, a mesma \emph{thread} atende outros eventos. Já o \textbf{cálculo pesado precisa de um núcleo} de verdade, fora da \emph{thread} da interface.
      \end{conceito}
    \end{column}
  \end{columns}
\end{frame}

\begin{frame}{O laço de eventos em cada plataforma}
  \begin{columns}[T]
    \begin{column}{0.55\textwidth}
      \begin{definicao}[Laço de eventos (\emph{event loop})]
        Uma \emph{thread} que repete sem parar: pega o próximo evento da fila, executa o código associado até o fim e volta à fila. Toda interface gráfica funciona assim.
      \end{definicao}
      \scriptsize
      \begin{tabularx}{\linewidth}{@{}l L@{}}
        \toprule
        \cabecalho{Plataforma} & \cabecalho{Mecanismo na \emph{thread} principal} \\
        \midrule
        Android & \cod{Looper} + \cod{MessageQueue} + \cod{Handler} \\
        iOS & \cod{RunLoop} e fila principal do GCD \\
        Dart & fila de microtarefas e fila de eventos \\
        Web & laço de eventos do navegador (tarefas e microtarefas) \\
        \bottomrule
      \end{tabularx}
    \end{column}
    \begin{column}{0.43\textwidth}
      \begin{pordentro}[Dart: duas filas]
        Antes de cada evento, o Dart esvazia a fila de \textbf{microtarefas} (continuações curtas de \cod{Future}s). Desde o Flutter 3.29, no Android e no iOS, o código Dart roda na própria \emph{thread} principal do app.
      \end{pordentro}
      \begin{atencao}[Um evento lento atrasa todos]
        Quadros se perdem (\emph{jank}) e, no Android, sem resposta a um toque por cerca de 5\,s aparece o diálogo ANR.
      \end{atencao}
    \end{column}
  \end{columns}
\end{frame}

\begin{frame}{Da pirâmide de \emph{callbacks} ao \cod{async}/\cod{await}}
  \centering
  \begin{tikzpicture}[every node/.style={font=\tiny}]
    \draw[seta] (0,0) -- (13.4,0);
    \foreach \x/\t/\d in {1.2/\emph{callbacks}/{funções chamadas\\ no fim da tarefa},4.4/\emph{Futures} e \emph{Promises}/{objeto que representa\\ o resultado futuro},7.8/{\cod{async}/\cod{await}}/{C\# (2012), Dart (2015),\\ JavaScript (2017)},11.4/concorrência estruturada/{Kotlin (2018):\\ escopos e cancelamento}}{
      \fill[ambar] (\x,0) circle (2pt);
      \node[font=\scriptsize\bfseries,text=azulpetroleo,above=1mm] at (\x,0) {\t};
      \node[rotulo,below=1mm] at (\x,0) {\d};
    }
  \end{tikzpicture}
  \raggedright
  \begin{columns}[T]
    \begin{column}{0.49\textwidth}
      \begin{definicao}[\emph{Future} (ou \emph{Promise})]
        Objeto que representa um resultado que ainda não existe. Começa \textbf{pendente} e termina \textbf{com valor} ou \textbf{com erro}. Quem o recebe pode esperar sem bloquear a \emph{thread}.
      \end{definicao}
      \centering
      \begin{tikzpicture}[x=0.5cm,y=0.45cm,every node/.style={font=\tiny}]
        \node[rotulo,anchor=east] at (-0.2,0.4) {\emph{thread}};
        \fill[petroleo!80] (0,0) rectangle ++(2,0.8); \node[white] at (1,0.4) {trecho 1};
        \fill[cinzaclaro] (2,0) rectangle ++(3,0.8); \node at (3.5,0.4) {outros eventos};
        \fill[petroleo!80] (5,0) rectangle ++(2,0.8); \node[white] at (6,0.4) {trecho 2};
        \fill[cinzaclaro] (7,0) rectangle ++(2.4,0.8); \node at (8.2,0.4) {outros};
        \fill[petroleo!80] (9.4,0) rectangle ++(1.6,0.8); \node[white] at (10.2,0.4) {fim};
        \foreach \x in {2,7}{\draw[ambar,line width=1pt] (\x,-0.2) -- (\x,1.1); \node[rotulo,above] at (\x,1.1) {\cod{await}};}
      \end{tikzpicture}
    \end{column}
    \begin{column}{0.49\textwidth}
      \begin{pordentro}[O que acontece em um \cod{await}]
        A função pausa em um \textbf{ponto de suspensão} e devolve a \emph{thread}; quando o resultado chega, continua de onde parou, com as variáveis locais preservadas. O compilador do Kotlin transforma cada função \cod{suspend} em uma máquina de estados que recebe uma \cod{Continuation}.
      \end{pordentro}
      \begin{dica}
        O ganho é escrever código assíncrono com cara de sequencial: \cod{try}/\cod{catch} e laços funcionam como sempre.
      \end{dica}
    \end{column}
  \end{columns}
\end{frame}

\begin{frame}{Concorrência estruturada}
  \begin{columns}[c]
    \begin{column}{0.45\textwidth}
      \centering
      \begin{tikzpicture}[every node/.style={font=\tiny}]
        \node[caixaescura,minimum width=3.4cm,font=\tiny] (e) at (0,0) {escopo da tela};
        \node[caixa,minimum width=1.6cm,font=\tiny] (a) at (-0.95,-1) {carregar lista};
        \node[caixa,minimum width=1.6cm,font=\tiny] (b) at (0.95,-1) {ouvir sensor};
        \node[caixaambar,minimum width=1.6cm,font=\tiny] (c) at (0.95,-2) {ler histórico};
        \draw[seta] (e) -- (a); \draw[seta] (e) -- (b); \draw[seta] (b) -- (c);
        \draw[pontilhado] (-2,0.5) rectangle (2,-2.5);
        \node[rotulo,text width=4.2cm] at (0,-3) {a tela fecha $\Rightarrow$ o escopo termina $\Rightarrow$ as tarefas filhas são canceladas};
      \end{tikzpicture}
      \begin{historia}
        Assim como a programação estruturada aposentou o \cod{goto}, este modelo aposenta tarefas soltas. A ideia ganhou força em 2018, com um ensaio de Nathaniel J. Smith, e o Kotlin a adotou no mesmo ano.
      \end{historia}
    \end{column}
    \begin{column}{0.53\textwidth}
      \begin{definicao}[Concorrência estruturada]
        Toda tarefa concorrente nasce dentro de um \textbf{escopo} e não vive mais que ele: quando o escopo termina, as tarefas filhas são canceladas.
      \end{definicao}
      \begin{conceito}[Garantias]
        Nada fica rodando depois que a tela fecha; o cancelamento chega aos filhos; erros dos filhos chegam ao pai.
      \end{conceito}
      \begin{atencao}[E no Dart?]
        \cod{Future} e \cod{Stream} não seguem esse modelo: cancele inscrições e temporizadores no \cod{dispose} e confira \cod{mounted} após um \cod{await}.
      \end{atencao}
    \end{column}
  \end{columns}
\end{frame}

\begin{frame}{Condições de corrida e impasses}
  \begin{columns}[T]
    \begin{column}{0.47\textwidth}
      \centering
      \begin{tikzpicture}[x=0.7cm,y=0.6cm,every node/.style={font=\tiny}]
        \tikzset{passo/.style={caixa,minimum height=4.5mm,minimum width=1.2cm,inner sep=2pt,font=\tiny}}
        \node[rotulo,anchor=east] at (-0.1,1) {\emph{thread} A};
        \node[rotulo,anchor=east] at (-0.1,0) {\emph{thread} B};
        \node[passo] at (0.9,1) {lê 5};
        \node[passo] at (2.6,0) {lê 5};
        \node[passo,caixaambar,font=\tiny] at (4.3,1) {grava 6};
        \node[passo,caixaambar,font=\tiny] at (6,0) {grava 6};
        \draw[->,cinzamedio] (0,-0.7) -- (7,-0.7) node[rotulo,right]{t};
        \node[rotulo,text=vermelhoerro,text width=4.8cm] at (3.4,-1.3) {duas somas de 1 a partir de 5 deveriam dar 7: uma atualização se perdeu};
      \end{tikzpicture}
      \begin{definicao}[Condição de corrida]
        Erro em que o resultado depende da ordem imprevisível em que tarefas concorrentes acessam um dado compartilhado.
      \end{definicao}
      \begin{atencao}[Impasse (\emph{deadlock})]
        Duas tarefas, cada uma segurando um recurso e esperando o da outra: ninguém avança. Adquira travas sempre na mesma ordem.
      \end{atencao}
    \end{column}
    \begin{column}{0.51\textwidth}
      \scriptsize
      \begin{tabularx}{\linewidth}{@{}l L@{}}
        \toprule
        \cabecalho{Proteção} & \cabecalho{Como funciona} \\
        \midrule
        Exclusão mútua & uma tarefa por vez na região crítica (\cod{Mutex}) \\
        Operação atômica & ler e gravar em um passo indivisível (\cod{AtomicInteger}, \cod{StateFlow.update}) \\
        Confinamento & só uma \emph{thread} toca o dado (a principal, por exemplo) \\
        Imutabilidade & ninguém altera: cria-se uma cópia nova \\
        Mensagens & cada tarefa tem seus dados e troca cópias \\
        \bottomrule
      \end{tabularx}
      \normalsize
      \begin{pordentro}[Corrida lógica com uma \emph{thread} só]
        Entre dois \cod{await}, outro evento pode mudar o estado. Duas buscas disparadas em sequência podem responder fora de ordem, e a mais antiga sobrescrever a mais nova.
      \end{pordentro}
    \end{column}
  \end{columns}
\end{frame}

\begin{frame}{Memória compartilhada ou troca de mensagens}
  \begin{columns}[c]
    \begin{column}{0.43\textwidth}
      \centering
      \begin{tikzpicture}[every node/.style={font=\tiny}]
        \node[caixa,minimum width=1.3cm,font=\tiny] (a) at (0,0) {\emph{thread} A};
        \node[caixa,minimum width=1.3cm,font=\tiny] (b) at (2.6,0) {\emph{thread} B};
        \node[caixaambar,minimum width=1.6cm,font=\tiny] (o) at (1.3,-1.2) {objeto + trava};
        \draw[seta] (a) -- (o); \draw[seta] (b) -- (o);
        \node[rotulo] at (1.3,0.6) {\textbf{memória compartilhada}};
        \node[caixaescura,minimum width=1.3cm,minimum height=9mm,font=\tiny] (c) at (0,-3) {isolate A\\ memória A};
        \node[caixaescura,minimum width=1.3cm,minimum height=9mm,font=\tiny] (d) at (2.6,-3) {isolate B\\ memória B};
        \draw[setaambar] ([yshift=1.5mm]c.east) -- node[rotulo,above]{mensagem} ([yshift=1.5mm]d.west);
        \draw[setaambar] ([yshift=-1.5mm]d.west) -- node[rotulo,below]{resposta} ([yshift=-1.5mm]c.east);
        \node[rotulo] at (1.3,-2.1) {\textbf{troca de mensagens}};
      \end{tikzpicture}
    \end{column}
    \begin{column}{0.55\textwidth}
      \scriptsize
      \begin{tabularx}{\linewidth}{@{}l L L@{}}
        \toprule
         & \cabecalho{Memória compartilhada} & \cabecalho{Troca de mensagens} \\
        \midrule
        Exemplos & \emph{threads} da JVM (Kotlin, Java) & \emph{isolates} do Dart, atores, Web Workers \\
        Comunicação & ler e gravar o mesmo objeto & enviar cópias de dados \\
        Risco & corridas e impasses & custo de copiar dados grandes \\
        Proteção & travas, atômicos, confinamento & o próprio modelo isola \\
        \bottomrule
      \end{tabularx}
      \normalsize
      \begin{analogia}
        Um quadro branco compartilhado exige regras para ninguém apagar o do colega. Colegas que trocam bilhetes nunca brigam pelo quadro, mas gastam tempo escrevendo.
      \end{analogia}
    \end{column}
  \end{columns}
\end{frame}

\begin{frame}{Fluxos de dados: frios, quentes e contrapressão}
  \begin{columns}[T]
    \begin{column}{0.55\textwidth}
      \begin{definicao}[Fluxo (\emph{stream})]
        Sequência de valores que chegam ao longo do tempo: leituras de sensor, mensagens, mudanças no banco. Segue o padrão \textbf{Observador}: quem produz notifica quem se inscreveu.
      \end{definicao}
      \scriptsize
      \begin{tabularx}{\linewidth}{@{}l L l l@{}}
        \toprule
        \cabecalho{Tipo} & \cabecalho{Comportamento} & \cabecalho{Kotlin} & \cabecalho{Dart} \\
        \midrule
        Frio & o produtor começa a cada inscrição & \cod{Flow} & \cod{Stream} (\cod{async*}) \\
        Quente & existe sem observadores; quem chega tarde perde o passado & \cod{SharedFlow} & \cod{Stream} \emph{broadcast} \\
        Estado & quente e guarda o valor atual & \cod{StateFlow} & \cod{ValueNotifier} \\
        \bottomrule
      \end{tabularx}
    \end{column}
    \begin{column}{0.43\textwidth}
      \begin{definicao}[Contrapressão (\emph{backpressure})]
        O que fazer quando o produtor é mais rápido que o consumidor: acumular (\emph{buffer}), descartar, manter só o mais recente (\emph{conflate}) ou amostrar (\emph{debounce}, \emph{sample}).
      \end{definicao}
      \begin{exemplo}[No LabMonitor]
        Um sensor envia 100 leituras por segundo e a tela desenha 60 quadros: basta mostrar sempre a leitura mais recente.
      \end{exemplo}
    \end{column}
  \end{columns}
\end{frame}


% =============================================================================
\section{O event loop}

\begin{frame}{Uma única thread, sem travar}
  \begin{columns}[c]
    \begin{column}{0.5\textwidth}
      \small
      \begin{itemize}
        \item O código Dart do app roda em um \termo{isolate} principal, com uma única thread.
        \item O \termo{event loop} processa, um de cada vez, toques, quadros a desenhar, respostas de rede e temporizadores.
        \item Para 60\,Hz, cada quadro tem cerca de \textbf{16\,ms}. Uma tarefa longa atrasa os quadros e a tela ``engasga'' (\emph{jank}).
        \item Operações de I/O (rede, arquivo, banco) são \textbf{assíncronas}: o Dart pede ao sistema e segue atendendo eventos até a resposta chegar.
      \end{itemize}
    \end{column}
    \begin{column}{0.48\textwidth}
      \centering
      \begin{tikzpicture}[node distance=3mm]
        \node[caixa,minimum width=1.4cm] (t) {toque};
        \node[caixa,minimum width=1.4cm,right=of t] (q) {quadro};
        \node[caixa,minimum width=1.4cm,right=of q] (r) {resposta HTTP};
        \node[pontilhado,fit=(t)(q)(r),label={[rotulo]above:fila de eventos}] (fila) {};
        \node[caixaescura,minimum width=3.2cm,below=10mm of q] (loop) {\emph{event loop}};
        \draw[seta] (fila) -- node[rotulo,right]{um por vez} (loop);
        \draw[setaambar,rounded corners] (loop.west) -- ++(-6mm,0) |- (fila.west);
      \end{tikzpicture}
      \begin{atencao}
        I/O não trava a UI se for assíncrono. Já um laço de CPU pesado trava, mesmo dentro de uma função \cod{async}: para isso existem os \emph{isolates}.
      \end{atencao}
    \end{column}
  \end{columns}
\end{frame}

% =============================================================================
\section{Future e async/await}

\begin{frame}[fragile]{\cod{Future}, \cod{async} e \cod{await}}
  \begin{columns}[T]
    \begin{column}{0.52\textwidth}
\begin{codigo}[title=future.dart]{Dart}
Future<double> lerTemperatura() async {
  await Future.delayed(const Duration(milliseconds: 500));
  return 24.7;
}

Future<String> relatorio() async {
  final t = await lerTemperatura();   // espera sem travar
  return 'Temperatura: $t °C';
}

void main() async {
  print('antes');
  final r = await relatorio();
  print(r);
  print('depois');
}
\end{codigo}
    \end{column}
    \begin{column}{0.46\textwidth}
      \begin{conceito}[\cod{Future<T>}]
        Promessa de um valor do tipo \cod{T} (ou de um erro) que chegará no futuro. Uma função \cod{async} sempre devolve um \cod{Future}.
      \end{conceito}
      \begin{conceito}[\cod{await}]
        Suspende a função até o \cod{Future} terminar, liberando o \emph{event loop} para outros eventos. Só pode ser usado dentro de funções \cod{async}.
      \end{conceito}
      \begin{erro}
        Esquecer o \cod{await} faz o código seguir sem o resultado. O \emph{lint} \cod{unawaited\_futures} ajuda a encontrar esses casos.
      \end{erro}
    \end{column}
  \end{columns}
\end{frame}

\begin{frame}[fragile]{Erros, tempo limite e paralelismo}
  \begin{columns}[T]
    \begin{column}{0.52\textwidth}
\begin{codigo}[title=Tratando falhas]{Dart}
Future<void> carregar() async {
  try {
    final dados = await api.bancadas()
        .timeout(const Duration(seconds: 10));
    mostrar(dados);
  } on TimeoutException {
    mostrarErro('O servidor demorou demais');
  } on SocketException {
    mostrarErro('Sem conexão');
  } catch (e, pilha) {
    registrar(e, pilha);
  }
}
\end{codigo}
\begin{codigo}[title=Em paralelo]{Dart}
final (temp, umid) = await (
  lerTemperatura(), lerUmidade()
).wait;             // os dois rodam ao mesmo tempo
\end{codigo}
    \end{column}
    \begin{column}{0.46\textwidth}
      \scriptsize
      \begin{tabularx}{\linewidth}{@{}l L@{}}
        \toprule
        \cabecalho{Recurso} & \cabecalho{Uso} \\
        \midrule
        \cod{.timeout(d)} & falha após um tempo \\
        \cod{Future.wait([..])} & espera vários \\
        \cod{(a, b).wait} & idem, com \emph{records} \\
        \cod{Future.any} & o primeiro que terminar \\
        \cod{Future.delayed} & espera um tempo \\
        \bottomrule
      \end{tabularx}
      \begin{dica}
        Sequencial: tempo total = soma. Em paralelo: tempo total $\approx$ a mais lenta.
      \end{dica}
      \begin{atencao}[Cancelamento]
        \cod{Future} não se cancela: use \cod{Stream} ou o \cod{CancelToken} do Dio (Módulo 11).
      \end{atencao}
    \end{column}
  \end{columns}
\end{frame}

% =============================================================================
\section{Streams}

\begin{frame}[fragile]{\cod{Stream}: vários valores no tempo}
  \begin{columns}[T]
    \begin{column}{0.52\textwidth}
\begin{codigo}[title=Simulando um sensor]{Dart}
Stream<double> temperaturas() async* {
  var t = 24.0;
  final rnd = Random();
  while (true) {
    t += rnd.nextDouble() * 0.6 - 0.3;
    yield t;                          // emite um valor
    await Future.delayed(const Duration(seconds: 1));
  }
}

final inscricao = temperaturas()
    .map((t) => t.toStringAsFixed(1))
    .listen((texto) => print('$texto °C'));

// ao sair da tela:
await inscricao.cancel();
\end{codigo}
    \end{column}
    \begin{column}{0.46\textwidth}
      \centering
      \begin{tikzpicture}[node distance=4mm]
        \node[caixaescura] (p) {produtor\\ \cod{yield}};
        \node[caixaambar,right=of p] (o) {operadores\\ \cod{map}, \cod{where}};
        \node[caixa,right=of o] (c) {ouvinte\\ \cod{listen}};
        \draw[seta] (p) -- (o); \draw[seta] (o) -- (c);
      \end{tikzpicture}
      \begin{conceito}[Future \emph{versus} Stream]
        \cod{Future}: \textbf{um} valor. \cod{Stream}: \textbf{vários} valores ao longo do tempo (leituras de sensor, mensagens MQTT, mudanças no banco).
      \end{conceito}
      \begin{atencao}
        Sempre cancele a inscrição (\cod{cancel()}) quando não precisar mais dos dados; senão o sensor continua gastando bateria.
      \end{atencao}
    \end{column}
  \end{columns}
\end{frame}

\begin{frame}[fragile]{Operadores e StreamController}
  \begin{columns}[T]
    \begin{column}{0.52\textwidth}
\begin{codigo}[title=Pipeline reativo]{Dart}
final alertas = temperaturas()
    .map((t) => double.parse(t.toStringAsFixed(1)))
    .distinct()                     // ignora repetidos
    .where((t) => t > 30)
    .map((t) => 'Alerta: $t °C');

await for (final a in alertas.take(3)) {
  print(a);                         // consome 3 e para
}
\end{codigo}
\begin{codigo}[title=StreamController]{Dart}
class Barramento {
  final _ctrl = StreamController<String>.broadcast();
  Stream<String> get eventos => _ctrl.stream;
  void publicar(String e) => _ctrl.add(e);
  void fechar() => _ctrl.close();
}
\end{codigo}
    \end{column}
    \begin{column}{0.46\textwidth}
      \scriptsize
      \begin{tabularx}{\linewidth}{@{}l L@{}}
        \toprule
        \cabecalho{Operador} & \cabecalho{Efeito} \\
        \midrule
        \cod{map}, \cod{where} & transformar e filtrar \\
        \cod{distinct} & descarta repetidos seguidos \\
        \cod{take}, \cod{first} & limita a quantidade \\
        \cod{asyncMap} & transforma com \cod{Future} \\
        \cod{handleError} & trata erros \\
        \cod{timeout} & erro se ficar sem eventos \\
        \bottomrule
      \end{tabularx}
      \begin{conceito}[Simples \emph{versus} broadcast]
        Uma \cod{Stream} simples aceita \textbf{um} ouvinte; uma \emph{broadcast} aceita vários.
      \end{conceito}
      \begin{dica}
        O pacote \cod{rxdart} acrescenta \cod{debounceTime}, \cod{combineLatest} e outros operadores.
      \end{dica}
    \end{column}
  \end{columns}
\end{frame}

% =============================================================================
\section{Assíncrono na interface}

\begin{frame}[fragile]{FutureBuilder e StreamBuilder}
  \begin{columns}[T]
    \begin{column}{0.56\textwidth}
\begin{codigo}[title=Painel em tempo real]{Dart}
class PainelTemperatura extends StatelessWidget {
  const PainelTemperatura({super.key, required this.leituras});
  final Stream<double> leituras;

  @override
  Widget build(BuildContext context) => StreamBuilder<double>(
    stream: leituras,
    builder: (context, snap) {
      if (snap.hasError) return Text('Erro: ${snap.error}');
      if (!snap.hasData) {
        return const CircularProgressIndicator();
      }
      return Text('${snap.data!.toStringAsFixed(1)} °C',
          style: Theme.of(context).textTheme.displayMedium);
    },
  );
}
\end{codigo}
    \end{column}
    \begin{column}{0.42\textwidth}
      \begin{conceito}[\cod{AsyncSnapshot}]
        \cod{snap.connectionState}, \cod{hasData}, \cod{data}, \cod{hasError} e \cod{error} permitem desenhar os estados carregando, sucesso e erro.
      \end{conceito}
      \begin{erro}[Future criado no \cod{build}]
        \cod{FutureBuilder(future: api.carregar())} refaz a requisição a cada reconstrução. Crie o \cod{Future} no \cod{initState} ou no \emph{view model} (Módulo 08).
      \end{erro}
    \end{column}
  \end{columns}
\end{frame}

\begin{frame}[fragile]{Temporizadores e \cod{mounted}}
  \begin{columns}[T]
    \begin{column}{0.54\textwidth}
\begin{codigo}[title=Atualização periódica]{Dart}
class _PainelState extends State<Painel> {
  Timer? _timer;
  double? _temp;

  @override
  void initState() {
    super.initState();
    _timer = Timer.periodic(const Duration(seconds: 5),
        (_) => _atualizar());
  }

  Future<void> _atualizar() async {
    final t = await api.temperatura();
    if (!mounted) return;            // a tela ainda existe?
    setState(() => _temp = t);
  }

  @override
  void dispose() { _timer?.cancel(); super.dispose(); }
  // build(...)
}
\end{codigo}
    \end{column}
    \begin{column}{0.44\textwidth}
      \begin{atencao}[Depois de um \cod{await}]
        O usuário pode ter saído da tela. Verifique \cod{mounted} antes de chamar \cod{setState} ou usar o \cod{context}.
      \end{atencao}
      \begin{conceito}[Limpeza]
        Tudo que é criado no \cod{initState} (timers, inscrições, controladores) é liberado no \cod{dispose}.
      \end{conceito}
    \end{column}
  \end{columns}
\end{frame}

% =============================================================================
\section{Isolates}

\begin{frame}[fragile]{Isolates: CPU sem travar a interface}
  \begin{columns}[T]
    \begin{column}{0.52\textwidth}
\begin{codigo}[title=Cálculo pesado em outro isolate]{Dart}
// função pura, sem acesso à UI
List<double> fft(List<double> amostras) {
  // ... milhares de operações
  return espectro;
}

Future<void> analisar() async {
  final espectro = await Isolate.run(
      () => fft(amostras));     // roda em paralelo
  setState(() => _espectro = espectro);
}

// atalho do Flutter (também funciona na web)
final dados = await compute(decodificarJson, textoGrande);
\end{codigo}
    \end{column}
    \begin{column}{0.46\textwidth}
      \centering
      \begin{tikzpicture}[node distance=4mm]
        \node[caixaescura,minimum width=2.4cm] (ui) {isolate principal\\ \tiny UI e eventos};
        \node[caixaambar,minimum width=2.4cm,right=12mm of ui] (w) {novo isolate\\ \tiny CPU};
        \draw[seta] ([yshift=2mm]ui.east) -- node[rotulo,above]{dados} ([yshift=2mm]w.west);
        \draw[setaambar] ([yshift=-2mm]w.west) -- node[rotulo,below]{resultado} ([yshift=-2mm]ui.east);
      \end{tikzpicture}
      \begin{conceito}[Isolate]
        Unidade de execução com \textbf{memória própria}: não compartilha variáveis, troca mensagens. Por isso não há condições de corrida.
      \end{conceito}
      \begin{dica}[Quando usar]
        FFT de sinais, filtros, compressão de imagem, JSON muito grande (centenas de KB) e criptografia.
      \end{dica}
    \end{column}
  \end{columns}
\end{frame}

% =============================================================================
\section{Fechamento}

\begin{frame}{Resumo}
  \begin{resumo}
    \begin{itemize}
      \item O Dart roda o app num \emph{event loop} de uma thread; nada pode bloqueá-lo por muito tempo.
      \item \cod{Future} representa um valor futuro; \cod{async}/\cod{await} deixam o código sequencial e legível.
      \item Trate erros com \cod{try}/\cod{on}/\cod{catch}, use \cod{timeout} e paralelize com \cod{Future.wait}.
      \item \cod{Stream} emite vários valores; cancele as inscrições quando terminar.
      \item \cod{FutureBuilder}/\cod{StreamBuilder} desenham os estados; verifique \cod{mounted} após \cod{await}.
      \item Cálculo pesado vai para outro isolate com \cod{Isolate.run} ou \cod{compute}.
    \end{itemize}
  \end{resumo}
\end{frame}

\begin{frame}{Glossário (1/2)}
  \begin{glossario}
    \emph{Thread} & linha de execução dentro de um processo, com pilha própria \\
    \emph{Isolate} & unidade de execução do Dart com memória própria; comunica-se por mensagens \\
    Laço de eventos & processa os eventos da fila, um de cada vez, na mesma \emph{thread} \\
    Microtarefa & tarefa curta executada antes do próximo evento da fila \\
    Concorrência & tarefas em andamento no mesmo período, que se alternam \\
    Paralelismo & tarefas executando no mesmo instante, em núcleos diferentes \\
    \cod{Future} & resultado que ainda vai existir: pendente, com valor ou com erro \\
    \cod{async}/\cod{await} & escrevem código assíncrono com cara de sequencial \\
    \emph{Jank} & quadro perdido porque a \emph{thread} da interface ficou ocupada \\
  \end{glossario}
\end{frame}

\begin{frame}{Glossário (2/2)}
  \begin{glossario}
    \cod{Stream} & fluxo de vários valores ao longo do tempo \\
    \cod{StreamSubscription} & inscrição em um fluxo; deve ser cancelada ao terminar \\
    \emph{Broadcast} & fluxo quente que aceita vários ouvintes \\
    \cod{StreamController} & cria e alimenta um fluxo manualmente \\
    Condição de corrida & resultado que depende da ordem imprevisível dos eventos \\
    Contrapressão & estratégia para quando o produtor é mais rápido que o consumidor \\
    \cod{Isolate.run} & executa uma função em outro \emph{isolate} e devolve o resultado \\
  \end{glossario}
\end{frame}

\begin{frame}{Questões de revisão}
  \begin{questoes}
    \item Se cada \emph{isolate} tem uma única \emph{thread}, como o app faz uma requisição de rede sem travar a tela?
    \item Diferencie concorrência de paralelismo. Qual deles exige outro \emph{isolate} no Dart?
    \item Por que um laço de CPU pesado trava a interface mesmo dentro de uma função \cod{async}?
    \item O que acontece no laço de eventos quando a execução chega a um \cod{await}?
    \item Quem roda primeiro: microtarefas ou eventos? Por que isso importa?
    \item Dê um exemplo de condição de corrida em Dart, mesmo com uma \emph{thread} só, e como evitá-la.
    \item Qual a diferença entre um \cod{Stream} de inscrição única e um \emph{broadcast}?
    \item Por que \emph{isolates} não precisam de travas para proteger seus dados?
  \end{questoes}
\end{frame}

\begin{frame}{Exercícios}
  \begin{exercicio}[Exercício 1 --- Paralelismo]
    Escreva três funções que simulam leituras com atrasos de 1\,s, 2\,s e 3\,s. Meça o tempo total em sequência e com \cod{Future.wait}.
  \end{exercicio}
  \begin{exercicio}[Exercício 2 --- Stream]
    Crie uma \cod{Stream<int>} que simula um ADC de 12 bits a 10\,Hz e um pipeline que converte para volts, agrupa 10 amostras, calcula a média e emite alerta acima de 3,0\,V.
  \end{exercicio}
  \begin{exercicio}[Exercício 3 --- Isolate]
    Gere 1 milhão de amostras e calcule estatísticas na thread principal com uma animação girando na tela. Repita com \cod{Isolate.run} e compare.
  \end{exercicio}
  \begin{projeto}[LabMonitor --- Etapa 6]
    Implemente \cod{SimuladorSensor} com uma \cod{Stream<Leitura>} por bancada e mostre a última leitura de cada bancada com \cod{StreamBuilder}.
  \end{projeto}
\end{frame}

\begin{frame}{Para aprofundar}
  \small
  \begin{itemize}
    \item Programação assíncrona: \url{https://dart.dev/libraries/async/async-await}
    \item Streams: \url{https://dart.dev/libraries/async/using-streams}
    \item Concorrência e isolates: \url{https://dart.dev/language/concurrency}
    \item Isolates no Flutter: \url{https://docs.flutter.dev/perf/isolates}
    \item Pacote rxdart: \url{https://pub.dev/packages/rxdart}
  \end{itemize}
\end{frame}

\begin{frame}{Respostas das questões de revisão}
  \begin{respostas}
    \item A requisição é I/O assíncrono: o Dart a entrega ao sistema operacional e segue atendendo eventos; quando a resposta chega, um evento completa o \cod{Future}.
    \item Concorrência: tarefas que se alternam no mesmo período (vários \cod{Future}s no laço). Paralelismo: execução simultânea em núcleos diferentes, que no Dart exige outro \emph{isolate}.
    \item \cod{async} não cria \emph{thread}: o trecho entre dois \cod{await} roda até o fim no laço de eventos. Um cálculo longo ocupa a \emph{thread} e impede os quadros; use \cod{Isolate.run}.
    \item A função pausa e devolve o controle ao laço, que atende outros eventos; quando o \cod{Future} completa, a continuação é agendada e a função segue do mesmo ponto.
    \item A fila de microtarefas é esvaziada antes do próximo evento. Microtarefas demais atrasam toques e quadros, como um laço longo.
    \item Duas buscas por texto: a resposta da primeira chega depois da segunda e sobrescreve o resultado mais novo. Guarde o identificador da última busca e ignore respostas antigas.
    \item Inscrição única aceita um só ouvinte e começa a produzir quando alguém escuta; o \emph{broadcast} aceita vários, e quem chega tarde perde os eventos anteriores.
    \item Cada \emph{isolate} tem memória própria e só troca cópias por mensagens: nenhum objeto mutável é acessado por dois \emph{isolates} ao mesmo tempo.
  \end{respostas}
\end{frame}

\encerramento{Módulo 07 --- Estado, navegação e animações: StatefulWidget, ciclo de vida, go\_router e animações.}

\end{document}
